% FORMALIZACION_NUEVA de la justificación del límite utilizado en
% 70_radiacion_termica.tex. No introduce un generador de constantes.
\subsection{Limite thermodynamique des sommes d’occupation}
\label{v:subsec:limite-termodinamico}

La densité radiale précédente peut être obtenue comme limite effective des
sommes de la réalisation périodique. L’estimation doit conserver le mode
constant fixé et contrôler uniformément aussi bien l’origine que la queue
spectrale. Les représentations d’action et de calendrier restent antérieures
à cette opération de réalisation.

\begin{lemma}[Convergence des trois observables thermiques]
\label{v:lem:limite-termodinamico}
Soient \(a,b>0\), \(L>0\), \(\Delta=2\pi/L\) et
\[
 F_N(r)=\frac{1}{e^{ar}-1},\qquad
 F_E(r)=\frac{br}{e^{ar}-1},\qquad
 F_Z(r)=-\log(1-e^{-ar})\quad(r>0).
\]
Pour chacune de ces fonctions, on a
\begin{equation}
 \lim_{L\to\infty}\frac{2}{L^3}
 \sum_{\nu\in\mathbb Z^3\setminus\{0\}}
 F\!\left(\frac{2\pi}{L}|\nu|\right)
 =\frac{1}{\pi^2}\int_0^\infty r^2F(r)\,dr.
 \label{v:eq:limite-modos-discretos}
\end{equation}
Les sommes et l’intégrale sont finies. La spécialisation
\(a=\beta\hbar c\), \(b=\hbar c\) fournit, respectivement,
les densités de nombre, d’énergie d’excitation et du logarithme de la fonction
de partition des modes transversaux.
\end{lemma}

\begin{proof}
Les trois fonctions sont continues et positives dans \((0,\infty)\).
Il existe des constantes \(C_0,C_1,d>0\), dépendant de \(a,b\) mais
indépendantes de \(\Delta\), telles que
\begin{equation}
 F(r)\leq\frac{C_0}{r}\quad(0<r\leq1),\qquad
 F(r)\leq C_1e^{-dr}\quad(r\geq1).
 \label{v:eq:envolventes-termicas}
\end{equation}
En effet, \(e^{ar}-1\geq ar\) contrôle \(F_N\) et \(F_E\)
près de l’origine ; en outre,
\(F_Z(r)=\log(1+F_N(r))\leq F_N(r)\leq1/(ar)\).
Pour \(r\geq1\), l’inégalité
\(-\log(1-y)\leq y/(1-y)\), avec \(y=e^{-ar}\), donne la borne
exponentielle de \(F_Z\). Les deux autres découlent de
\((e^{ar}-1)^{-1}\leq e^{-ar}/(1-e^{-a})\), en absorbant le facteur
\(r\) de \(F_E\) dans une exponentielle de taux inférieur. Ces bornes
prouvent déjà l’intégrabilité de \(r^2F(r)\).

Pour les sommes discrètes, nous regroupons les points entiers selon
\(m=\|\nu\|_\infty\geq1\). Chaque couche contient exactement
\[
 (2m+1)^3-(2m-1)^3=24m^2+2
\]
points et satisfait \(m\leq|\nu|\leq\sqrt3m\).
Fixons \(0<\delta\leq1\) et \(0<\Delta\leq1\).
Si \(\delta<\Delta\), il n’y a pas de points avec
\(0<\Delta|\nu|\leq\delta\). Sinon, en posant
\(N=\lfloor\delta/\Delta\rfloor\), la première borne de
\eqref{v:eq:envolventes-termicas} permet d’écrire
\begin{align*}
 \Delta^3\!\sum_{0<\Delta|\nu|\leq\delta}F(\Delta|\nu|)
 &\leq C_0\Delta^2\sum_{m=1}^{N}\frac{24m^2+2}{m}\\
 &\leq26C_0\Delta^2\sum_{m=1}^{N}m
 \leq26C_0\delta^2.
\end{align*}
Dans la deuxième ligne, on a utilisé \(m^{-1}\leq m\) puis
\(N(N+1)/2\leq N^2\) pour \(N\geq1\).
La contribution intégrale du même voisinage est aussi d’ordre
\(\delta^2\), car
\(\int_0^\delta r^2F(r)\,dr\leq C_0\delta^2/2\).

Pour \(R\geq1\), tout point avec \(\Delta|\nu|>R\) appartient
à une couche \(m>R/(\sqrt3\Delta)\). La seconde borne de
\eqref{v:eq:envolventes-termicas} donne
\begin{align*}
 \Delta^3\!\sum_{\Delta|\nu|>R}F(\Delta|\nu|)
 &\leq C_1e^{-dR/(2\sqrt3)}
 \Delta^3\sum_{m\geq1}(24m^2+2)e^{-d\Delta m/2}\\
 &\leq C_2e^{-dR/(2\sqrt3)},
\end{align*}
où \(C_2\) ne dépend pas de \(0<\Delta\leq1\). Pour vérifier
cette uniformité, on prend \(q=e^{-d\Delta/2}\) et on utilise
\[
 \sum_{m\geq1}q^m=\frac{q}{1-q},\qquad
 \sum_{m\geq1}m^2q^m=\frac{q(1+q)}{(1-q)^3}.
\]
La première identité provient de la série géométrique et la seconde de
l’application à deux reprises de \(q\,d/dq\), opération valable dans \(|q|<1\).
En outre,
\(1-e^{-d\Delta/2}\geq\Delta(1-e^{-d/2})\) par concavité.
Le facteur \(\Delta^3\) compense donc les dénominateurs des
deux expressions. La queue intégrale tend elle aussi vers zéro par la
borne exponentielle. Pour chaque \(\Delta>0\), ces mêmes estimations
prouvent la convergence de la somme infinie.

Dans la couronne compacte \(\delta\leq|k|\leq R\), la fonction
\(F(|k|)\) est continue et bornée. Les sommes de Riemann de maille
\(\Delta\) convergent vers son intégrale ; les deux sphères du bord
sont de mesure nulle. On fait d’abord \(\Delta\to0\), ensuite
\(\delta\to0\) et enfin \(R\to\infty\), en utilisant les bornes
uniformes précédentes. Il en résulte
\[
 \Delta^3\sum_{\nu\ne0}F(\Delta|\nu|)
 \longrightarrow \int_{\mathbb R^3}F(|k|)\,d^3k
 =4\pi\int_0^\infty r^2F(r)\,dr.
\]
Comme \(2/L^3=2\Delta^3/(2\pi)^3\), le coefficient radial final
est \(2\cdot4\pi/(2\pi)^3=1/\pi^2\), ce qui démontre
\eqref{v:eq:limite-modos-discretos}.
\end{proof}

Le contrôle de l’origine est effectué sur des excitations d’énergie positive.
Le mode constant reste fixé avant de former la fonction de partition :
son élimination ultérieure d’une intégrale ne remplacerait pas cette condition de
la représentation finie. La preuve établit le passage à la limite de la
réalisation déclarée et conserve la distinction entre cette réalisation et
la construction préalable de ses paramètres d’action et de calendrier.
